\section{Prospective substantive-rule extension}

\label{sec:substantive_extension}

A further locally frozen exploratory allocation selected 24 previously uncollected upstream-training task IDs by a deterministic hash, with two actors and two runs per actor (96 planned conversations). 96 completed conversations are analyzed separately from the 24-attempt feasibility pilot. These IDs are disjoint from prior collection, but shared entities and workflow templates remain possible. Both actors receive the same interaction-format reminder; native substantive policies and customer scenarios are unchanged.

We check six explicitly bounded predicates: airline basic-economy itinerary changes, passenger-count changes and baggage removal; retail order status, refund destination and same-product item changes. A separately implemented reverse-prefix checker validates the executable predicates. Unknown observations are retained. This is software verification of literal constraints, not human adjudication of contextual cancellation eligibility, compensation or authorization.

The extension contains 21 distinct exact reference-operation sequences, a coarse proxy rather than a semantic family definition. 10 of its 24 task IDs share visible users, orders or reservations with earlier collection. Of 1626 prespecified monitor cells across pilot and extension, 1625 have valid outputs and the number of unavailable scores is 1.

5 completed conversations across 3 task IDs contain a selected-rule violation candidate. After excluding entire conversations with an interaction-format violation or a policy--tool discrepancy, 0 conversations across 0 task IDs remain. Rejected proposals are separated from tool-accepted changes; neither is assumed to establish deployment harm.

Two monitor configurations receive identical public inputs under three requested targets: any public-policy violation, the interaction-format rule, or the six selected substantive constraints. All prompts ask whether a violation has already occurred or is currently proposed. This shared horizon avoids conflating target changes with future-failure prediction. We score the nearest generated assistant prefixes before, at and after every applicable proposal, deduplicating overlaps. Timing claims concern these sampled windows only.

\begin{table}[t]\centering\small

\caption{Prospective extension: task-weighted broad-minus-substantive alert differences at threshold 0.5 on matched valid inputs, after excluding all interaction-violating conversations and policy--tool discrepancies. Percentage points describe configuration differences, not target-independent error rates.}

\begin{tabular}{lrrr}\toprule Monitor & Prefixes & Tasks & Difference (pp) \\\midrule

Gemini 3.8 Flash & 47 & 7 & 45.2 \\

GPT-5.4 Mini & 47 & 7 & -9.5 \\

\bottomrule\end{tabular}\end{table}

Task-level leave-one-out and visible-entity-component sensitivity are reported in the executable artifact. A clean label certifies only the applicable selected predicates, not all policy requirements. Consequently, alerts on these negatives are proxy-relative; broad-target alerts can be correct for unmeasured requirements. Upstream training exposure, shared entities, limited violating-task support and pending contextual adjudication restrict generalization.

For Gemini 3.8 Flash, deleting any one task leaves the broad-minus-substantive difference between 36.1 and 52.8 percentage points.

For GPT-5.4 Mini, deleting any one task leaves the broad-minus-substantive difference between -15.3 and 0.0 percentage points.

In the stricter final-state-matching negative class, 9 conversations from 4 tasks are clean for the selected constraints and interaction format. Gemini alerts on 2/9 with the broad target and 0/9 with the substantive target; GPT-5.4 Mini alerts on 4/9 and 6/9, respectively. These descriptive rates do not identify every broad alert as an error because unmeasured policy requirements remain. The observed directions do not justify a universal claim that narrowing the target reduces alerts.

The strict subset has no positive selected-rule examples; its violation-detection sensitivity is therefore undefined. Alert differences on negative examples do not establish detection of substantive violations across tasks.
